% Prose draft of Background §2.3 (Shapes in Medical Image Analysis), moved here 2026-10-23
% when the section in text.tex was converted to bullet points. Not \input anywhere.

\section{Shapes in Medical Image Analysis} %7 pages
\label{sec:bg-shapes}

Medical image analysis uses a wide range of shape representations: binary masks, point
clouds, meshes, distance fields, statistical shape models, learned latent codes, and
neural implicit fields \citep{heimann2009statistical, ambellan2019statistical}. On first
encounter they look quite different, and the literature often treats them as separate
worlds. This section takes the opposite view: most of them are the same kind of object
seen through a few shared choices --- what one queries, what comes back, and how the
geometry is read off the answer. Making those choices explicit relates the
representations to one another and shows exactly where they differ. The aim is clarity of
classification rather than a survey: enough structure to say precisely how, for example,
a voxel mask and a neural network encoding the same organ are related.

\subsection{Representation as a queryable map}
\label{sec:shape-frame}

We fix a single object of interest. Let $S \subset \mathbb{R}^3$ denote the
\emph{shape} as a solid region --- the set of points the anatomy occupies --- and let
$\partial S$ denote its \emph{boundary surface}. These are the geometric ground truth;
everything below is a way of encoding one of them so that it can be recovered.

Formally, take a \emph{shape representation} to be a pair $(f, \mathcal{R})$. The first
component is a queryable map
\[
    f : \Omega \longrightarrow \mathcal{C},
\]
where $\Omega$ is the \emph{query domain} (the set of things one is allowed to ask
about) and $\mathcal{C}$ is the \emph{codomain} (the kind of answer that comes back).
The second component is a \emph{recovery operator} $\mathcal{R}$, a fixed rule that
reads the geometric shape --- $S$ or $\partial S$ --- off the map $f$. Representations
differ almost entirely in these three slots: \emph{what you query} ($\Omega$),
\emph{what you get back} ($\mathcal{C}$), and \emph{how geometry is recovered}
($\mathcal{R}$). The virtue of this view is that it exposes structural kinship: two
representations that agree on all three slots are, for our purposes, the same
representation stored differently.

Two recovery modes account for essentially all cases, and it is worth stating them
sharply because the entire explicit/implicit distinction rests on them.

\paragraph{Image recovery (explicit).} The stored values \emph{are} points or samples
of the shape itself. The codomain is ambient space, $\mathcal{C} = \mathbb{R}^3$, and
the recovered geometry is the image of $f$ --- possibly after a reconstruction step.
Here $\mathcal{R}$ is essentially ``read off the coordinates'': $f(\Omega)$ lies
\emph{on} $\partial S$ (or is $\partial S$). Crucially the definition is that the
stored values are points \emph{of} the shape, not the stronger claim that the shape
equals $f(\Omega)$; the latter holds for a mesh but not for a sparse landmark set or a
finite boundary sample, which are samples that must be labelled or reconstructed to
recover $S$.

\paragraph{Level-set recovery (implicit).} The query domain is space itself,
$\Omega \subseteq \mathbb{R}^3$, and $f$ returns a \emph{scalar} that is tested
against a threshold $\tau$. The shape is not stored; it is the sub- or superlevel set
of $f$. Depending on the sign convention one recovers the solid as
$S = f^{-1}\big((-\infty,\tau]\big)$ or $S = f^{-1}\big([\tau,\infty)\big)$, and the
boundary as the level set
\[
    \partial S = f^{-1}(\tau) = \{\, x \in \Omega : f(x) = \tau \,\}.
\]
The last equality is exact only under mild regularity (a non-vanishing gradient, or at
least no flat patch sitting exactly at value $\tau$); we flag this once and do not
repeat it. The sign convention must be stated with each field: for a signed distance
function, which is \emph{negative inside}, the solid is the sublevel set $\{d \le 0\}$;
for an occupancy or mask field, which is $\approx 1$ inside, the solid is the
\emph{superlevel} set $\{f \ge \tau\}$. Writing $\le$ for an occupancy field, or $\ge$
for an SDF, silently inverts the shape.

\paragraph{Classification, not distance.} This view is deliberately \emph{taxonomic}:
it tells us what a representation \emph{is} and how it relates to its neighbours. It is
not, by itself, a source of \emph{distance}. To measure how far two shapes lie apart one
first recovers a comparable object and then reaches for the appropriate metric --- and
the choice of object dictates the metric. Recovering a \emph{set} (a subset of
$\mathbb{R}^3$) invites the natural overlap and boundary comparisons: overlap scores such
as the Dice coefficient and boundary distances such as the (symmetric) Hausdorff
distance. Recovering a \emph{measure} (a mass distribution on $\mathbb{R}^3$) instead
opens the door to optimal-transport (Wasserstein) and information-theoretic (KL)
comparisons, which behave gracefully under blur and partial mass. The division of labour
is worth keeping explicit: the taxonomy \emph{classifies} representations, whereas a
metric is computed only \emph{after} a set or a measure has been recovered from one. We
treat the taxonomy as the organizing spine, because it is where the representational
choices live, and reach for a metric downstream whenever a distance is actually needed.

\subsection{The representations}
\label{sec:shape-reps}

We now walk through the canonical representations, each in its most standard form and
each marked by an icon in the margin. For every one we name the same slots --- what is
stored, what one queries ($\Omega$), what comes back ($\mathcal{C}$), and how the geometry
is recovered ($\mathcal{R}$) --- together with a plain intuition and the caveat most
easily got wrong.

\repicon{05_binary_mask.png}%
\paragraph{1.~Binary / voxel mask.} The most familiar representation: a segmentation
stored as a labelled raster \citep{isensee2021nnunet}. The domain is a \emph{finite grid}
$G \subset \mathbb{Z}^3$ (not all of ambient space), and the map $f : G \to \{0,1\}$ ---
or $[0,1]$ for a soft mask --- reports whether each voxel lies inside the object. Recovery
is by threshold: with inside $\approx 1$, the solid is the set of voxels above one half,
$S \approx \{\, x : f(x) \ge \tfrac12 \,\}$. Intuition: the organ painted onto a grid.
Caveat: whether a mask counts as \emph{implicit} is a matter of convention. By our
recovery test it is a thresholded field, hence implicit; but much of the neural-field
literature reserves ``implicit'' for continuous coordinate functions and calls a voxel
grid explicit and discrete. We classify it as implicit-by-recovery and simply flag the
tension.

\repicon{01_landmarks.png}%
\paragraph{2.~Landmarks / keypoints.} A small set of \emph{labelled} points, each label
naming one anatomically meaningful location: the domain is the index set
$\Omega = \{1, \dots, m\}$ and $f : \Omega \to \mathbb{R}^3$ gives each label a position
\citep{cootes1995active}. Recovery is explicit --- the stored values are points on or near
the surface $\partial S$. Intuition: named anatomical dots. Caveat: landmarks do
\emph{not} determine the whole shape; they are a sparse sample, valuable precisely because
their labels establish correspondence from one shape to the next.

\repicon{02_point_cloud.png}%
\paragraph{3.~Point cloud.} The same slots as landmarks --- a finite set of points in
$\mathbb{R}^3$ --- but now the points are \emph{unordered and unlabelled}: there is no
correspondence between one shape's points and another's
\citep{qi2017pointnet, wang2019dgcnn}. Recovery is explicit: the points are a sample
\emph{of} the boundary $\partial S$, and recovering the solid $S$ needs a
surface-reconstruction step \citep{kazhdan2006poisson}. Intuition: a scatter of surface
hits, as from a depth sensor or the surface voxels of a segmentation. The only formal
difference from landmarks is the missing labels --- and that single difference is what
makes point clouds cheap and abundant but harder to compare across shapes.

\repicon{03_mesh.png}%
\paragraph{4.~Mesh / surface.} A mesh stores a discrete object --- vertices
$V \subset \mathbb{R}^3$ joined into triangular faces $F$ --- but \emph{realizes} a
continuous one: interpolating linearly across each triangle yields a piecewise-flat
surface in space \citep{lorensen1987marching}. Recovery is explicit: the interpolated
triangles \emph{are} the boundary $\partial S$. Intuition: a triangulated skin stretched
over the organ. Caveat: the surface bounds a well-defined \emph{solid} only when the mesh
is watertight --- closed, consistently oriented, and free of self-intersections.
Anatomical meshes routinely violate this, and genuinely open surfaces (a cortical sheet, a
valve leaflet) bound no solid at all. The lesson is the stored-versus-realized split:
discrete data, continuous geometry, bridged by interpolation.

\repicon{04_sdf.png}%
\paragraph{5.~Distance function (SDF and UDF).} Instead of marking inside from outside,
store at every point in space its distance to the surface. The \emph{signed} distance
function (SDF) $d : \mathbb{R}^3 \to \mathbb{R}$ is negative inside and positive outside,
so recovery is again by threshold: the solid is $S = \{\, x : d(x) \le 0 \,\}$ and the
surface is the zero level set $\partial S = \{d = 0\}$
\citep{osher1988fronts, park2019deepsdf}. Intuition: at every point, how far to the
surface and on which side. It describes the same region as a mask, but as a \emph{smooth}
field, which is what makes it convenient for gradient-based optimization. Two facts are
worth keeping in mind. An exact SDF adds no information beyond the set $S$ --- its values
are fixed once $S$ is known; the payoff is computational (cheap distance and gradient
queries), and only \emph{learned} or approximate SDFs carry genuinely new information. And
away from the surface an SDF changes at unit rate, $|\nabla d| = 1$, except along the
object's skeleton, where the nearest surface point is ambiguous. The \emph{unsigned}
distance function (UDF) $u(x) = \operatorname{dist}(x, \partial S)$ drops the sign; it can
represent open surfaces that enclose no interior, at the price that no threshold brackets
the surface cleanly, so extracting it is harder \citep{chibane2020ndf}.

\repicon{06_ssm.png}%
\paragraph{6.~Statistical shape model (point-distribution form).} A statistical shape
model (SSM) represents not one shape but a whole \emph{population}, each member described
by $n$ corresponding points \citep{cootes1995active, heimann2009statistical}. Any shape is
written
\[
    s = \bar{s} + \Phi\,b,
\]
a mean configuration $\bar{s}$ plus a short coefficient vector $b$ that moves along a small
set of learned variation modes $\Phi$ (typically the leading PCA directions). Placing a
Gaussian on $b$ turns the model into a \emph{prior} over plausible anatomy
\citep{ambellan2019statistical}. Recovery is explicit: fixing $b$ produces a set of
landmarks. The point to carry forward is the split between a compact, shape-specific code
$b$ and a shared generator $(\bar{s}, \Phi)$ --- exactly the structure the learned latent
code below repeats. Caveat: this is the classical point-distribution form; SSMs can equally
be built on distance fields, in which case they are implicit.

\repicon{07_latent_code.png}%
\paragraph{7.~Latent code / embedding.} On its own a latent code represents
\emph{nothing}: a vector $z \in \mathbb{R}^d$ is just coordinates in an abstract space, with
no rule for turning it back into geometry. The shape is the \emph{pair} $(z, D_\theta)$ of
the code and a learned decoder $D_\theta$; geometry appears only through the field
$D_\theta(z, \cdot)$, and it is then classified by that field's recovery mode (usually
implicit --- an occupancy or SDF decoder) \citep{park2019deepsdf, mescheder2019occupancy}.
This mirrors the SSM directly: the code $z$ plays the role of the coefficient $b$, and the
decoder $D_\theta$ the role of the shared generator $(\bar{s}, \Phi)$. The difference is
that the SSM generator is linear and explicit, whereas the decoder is nonlinear and
implicit. Intuition: $z$ is the shape's compressed address; the decoder is the dictionary
that turns it back into geometry.

\repicon{08_inr.png}%
\paragraph{8.~Implicit neural representation (INR).} An INR stores the shape directly in the
weights $\theta$ of a small network $f_\theta : \mathbb{R}^3 \to \mathbb{R}$ that maps a
coordinate to a field value --- occupancy or signed distance
\citep{mescheder2019occupancy, park2019deepsdf, sitzmann2020siren}. Recovery is level-set,
exactly as for the distance function: threshold the network's output and read off the
surface. Intuition: rather than tabulate the field on a grid, a network learns it and can be
evaluated at any point and any resolution. The taxonomy makes the punchline plain: an INR
and a voxel-stored SDF have the \emph{same} codomain and the \emph{same} recovery rule. They
differ only in how the field is stored --- network weights versus a grid of samples --- and
in that the network is continuous where the grid is not. The representation that looks most
novel is, structurally, a continuously stored SDF.

\repicon{09_deformation.png}%
\paragraph{9.~Deformation field / diffeomorphism of a template.} One representation is special
to registration and central to this thesis: describe a shape not directly but as a
\emph{deformation of a reference}. A fixed template or atlas $S_0$ is warped by a spatial
transformation $\varphi : \mathbb{R}^3 \to \mathbb{R}^3$ --- often a diffeomorphism, a smooth
invertible warp --- so that the represented shape is $S = \varphi(S_0)$
\citep{beg2005lddmm, balakrishnan2019voxelmorph}. Both what one queries and what comes back are
ambient points: $\varphi$ maps space to space. Recovery is explicit --- apply $\varphi$ to the
template (or, if $S_0$ is stored as a field, warp the field by composition). Intuition: don't
describe the organ, describe how to bend a known organ into it. This is the precise sense in
which registration converts a comparison of \emph{shapes} into a comparison of
\emph{deformations}, and it is why this view sits at the centre of the thesis.

We set aside two further, more specialized families: spherical-harmonic (SPHARM)
parametrizations of closed surfaces \citep{styner2006framework}, and medial or
\mbox{$s$-rep} representations that encode a shape through its skeleton and a thickness field
\citep{pizer2003deformable}. Both matter, but neither is needed for what follows.

\begin{table}[htbp]
    \centering
    \footnotesize
    \setlength{\tabcolsep}{4pt}
    \renewcommand{\arraystretch}{1.25}
    \begin{tabularx}{\textwidth}{@{}l l l l X l l l@{}}
        \toprule
        Representation & $\Omega$ & $\mathcal{C}$ & Stored object & Recovery $\mathcal{R}$ & E/I & G/L & D/C \\
        \midrule
        Binary / voxel mask
            & $G \subset \mathbb{Z}^3$ & $\{0,1\}$ & occupancy raster
            & superlevel set $\{f \ge \tfrac12\}$ & I\textsuperscript{a} & G & D \\
        Landmarks / keypoints
            & $\{1,\dots,m\}$ & $\mathbb{R}^3$ & labelled points
            & read coordinates (sparse) & E & G & D \\
        Point cloud
            & $\mathbb{R}^{3n}/\mathfrak{S}_n$ & $\mathbb{R}^3$ & unlabelled points
            & sample of $\partial S$; reconstruct $S$ & E & G & D \\
        Mesh / surface
            & $|K|$ (realized) & $\mathbb{R}^3$ & vertices $V$, faces $F$
            & image $\varphi(|K|)$ & E & G & D\textsuperscript{b} \\
        Distance function (SDF)
            & $\mathbb{R}^3$ & $\mathbb{R}$ & scalar field $d$
            & sublevel set $\{d \le 0\}$ & I & G & C\textsuperscript{c} \\
        Distance function (UDF)
            & $\mathbb{R}^3$ & $\mathbb{R}_{\ge 0}$ & scalar field $u$
            & zero set $\{u=0\}$ (no interior) & I & G & C\textsuperscript{c} \\
        Statistical shape model
            & $\{1,\dots,n\}\!\times\!\mathbb{R}^k$ & $\mathbb{R}^3$ & $\bar{s}, \Phi$; measure on $\mathbb{R}^k$
            & per-vertex $\bar{s}_i + \Phi_i b$ & E\textsuperscript{d} & L & C \\
        Latent code / embedding
            & $\{*\}$ & $\mathbb{R}^d$ & vector $z$ (+ decoder)
            & none alone; via $D_\theta(z,\cdot)$ & I\textsuperscript{e} & L & C \\
        Implicit neural repr.\ (INR)
            & $\mathbb{R}^3$ & $\mathbb{R}$ & weights $\theta$
            & level set of $f_\theta$ & I & L & C \\
        Deformation field
            & $\mathbb{R}^3$ & $\mathbb{R}^3$ & warp $\varphi$ (+ template)
            & pushforward $\varphi(S_0)$ & E & G\textsuperscript{f} & C \\
        \bottomrule
    \end{tabularx}
    \caption[Taxonomy of shape representations]{Shape representations in the taxonomy of
    \secref{sec:shape-frame}. Columns give the query domain $\Omega$, codomain $\mathcal{C}$, the
    stored object, the recovery operator $\mathcal{R}$, and the three descriptive axes:
    Explicit/Implicit (E/I), Geometric/Learned (G/L), Discrete/Continuous (D/C).
    \textsuperscript{a}~Implicit \emph{by recovery mode}; a defensible choice, contested by
    literature that calls voxel grids explicit and discrete.
    \textsuperscript{b}~Discrete storage $(V,F)$; continuous realized surface $\varphi(|K|)$.
    \textsuperscript{c}~Continuous in concept, usually grid-sampled in practice.
    \textsuperscript{d}~Point-distribution form; level-set SSMs are implicit.
    \textsuperscript{e}~The code alone recovers nothing; the classification is inherited from the
    decoded field. \textsuperscript{f}~The warp is geometric; a network-predicted or
    velocity-parametrized warp is learned.}
    \label{tab:shape-representations}
\end{table}

\subsection{Three descriptive axes}
\label{sec:shape-axes}

The comparison in \tabref{tab:shape-representations} labels every representation along three axes.
It is tempting to say these axes ``fall out'' of the tuple $(\Omega, \mathcal{C}, f, \mathcal{R})$,
but that overstates the case. The taxonomy \emph{organizes} the axes as sharp questions one asks of the
slots and of information the slots do not fully contain; two of the three depend on facts about
provenance and storage that live outside the tuple.

\paragraph{Explicit vs.\ implicit --- a question about $\mathcal{R}$.} This one \emph{is} read off
the taxonomy, from the recovery operator alone. Do the stored values \emph{name} the geometry
(explicit: points of the shape, image recovery), or do they \emph{test} space against a threshold
(implicit: a scalar field, level-set recovery)? Landmarks, point clouds, meshes, SSMs, and
deformation fields are explicit; masks, distance functions, latent-code-plus-decoder, and INRs are
implicit.

\paragraph{Geometric vs.\ learned --- a question about provenance.} This is \emph{not} in the tuple. We
define \emph{learned} precisely: a representation is learned when its parameters $\theta$ are
\emph{optimization-fitted \emph{and} required at query time}. INR weights, a decoder $D_\theta$, and
the PCA basis $\Phi$ all qualify --- you cannot evaluate the representation without them, and they were
fitted. This definition has two deliberate consequences. It classifies even a \emph{single-shape} INR
as learned (its weights are fitted, though no population is involved), and it keeps masks, meshes, and
landmarks \emph{geometric} even when a trained network produced them, because the trained parameters are
not part of the representation at query time --- only the resulting coordinates or raster are. Note also
that ``learned'' is not the same as ``statistical''. Statistical means a probability measure lives on
$\Omega$ (as with the SSM's $\mathcal{N}(0,\Lambda)$); a single-shape INR is learned but not
statistical, and one can imagine geometric-but-statistical constructions. The axis is about fitted
parameters at query time, nothing more.

\paragraph{Discrete vs.\ continuous --- a stored-vs-realized question.} This too is not settled by the
tuple; it is a distinction between what is \emph{stored} and what is \emph{realized}. A mesh stores a
discrete $(V, F)$ but realizes a continuous surface through $\varphi$; a voxel-SDF tabulates a grid but
conceptually represents a continuous field; an INR stores continuous weights and realizes a continuous
field with no grid at all. The honest reading is to name both: masks and point sets are discrete in
storage, distance fields and neural fields are continuous in realization, and meshes sit astride the
two by construction.

Two representations can therefore share their answer on one axis and diverge on another, which is the
whole point of keeping the axes separate rather than collapsing them into a single ``modernness''
score. The sharpest example, again, is the INR against the voxel-SDF: identical on the
explicit/implicit axis (both implicit, same $\mathcal{R}$), identical in what they encode, and
separated \emph{only} on the discrete/continuous axis and by whether their parameters are fitted.

\subsection{What the taxonomy buys, and what it costs}
\label{sec:shape-closing}

The taxonomy earns its keep by making kinship and difference legible: it tells us that a latent code is a
generalized SSM coefficient, that an INR is a continuously stored SDF, that landmarks and point clouds
differ only in a label on their domain. What it does \emph{not} do is measure: it classifies structure,
and is taxonomic, not metric. The moment a distance is required --- overlap, boundary error,
transport cost --- one recovers a set or a measure from the representation and
computes there. Keeping this division explicit prevents a common confusion, namely reaching for a set
metric on an object (a bare latent code, a UDF) from which no set has yet been recovered.

Read in sequence, the representations trace an arc that this thesis will follow. It begins with direct,
geometric encodings --- masks, point clouds, meshes --- that store the anatomy as it is. It passes through
distance fields, which re-express the same region as a smooth, differentiable landscape and so make shape
amenable to gradient-based optimization. It ends with statistical and neural codes, which compress a shape
to a handful of numbers plus a shared generator and thereby capture what a \emph{population} of anatomies
have in common. The pay-off of this progression is exactly what makes shape useful for registration: a
representation abstract enough to be shared across images decouples the geometry of an organ from the
intensities of any one acquisition. That is what lets a shape serve as a common language --- a bridge ---
across differences in modality, scanner, and acquisition protocol, which is the theme the following
chapters develop.
